\documentclass[10pt]{article}
\usepackage[margin=0.76in]{geometry}
\usepackage[T1]{fontenc}
\usepackage{lmodern}
\usepackage{hyperref}
\title{Tactix: Sweep Gate and Learning-Readiness Checkpoint}
\author{Report to Andrew}
\date{25 September 2026}
\begin{document}
\maketitle
\vspace{-1.5em}
\emergencystretch=1em

\section*{Outcome}
The AM427 Sweep scheduling candidate remains unaccepted. Matched Tier A and B
are independently certificate-clean, but neither establishes N4 parity or a
stable performance gain. Tier C is still running on the same pinned binaries;
no Tier D comparison has been submitted. A source review confirms that
production online learning needs policy/session refactoring, while the existing
root-action executor supports bounded opt-in prototypes. The
\href{../../benchmarks/reports/2026-09-25-am429-sweep-tier-b-and-slurm-lock.md}{AM429}
and \href{../../benchmarks/reports/2026-09-25-am430-policy-lifetime-and-scheduling-fixture.md}{AM430}
technical reports hold the detailed evidence.

\section*{Verified evidence}
Tier A: both arms solved 199/200, including 117 independently checked native
DSR proofs each, with identical paired verdicts and proof sizes. AM427 took
3.0 more solver seconds across 200 tasks. Tier B: both arms solved 28/28
(14 SAT, 14 UNSAT), with all 14 proofs and all models checked; paired verdicts
and proof sizes again matched. AM427 took 415.978 versus the parent's 427.866
solver seconds over 28 tasks, a small, non-gating advantage. The first Tier C
submission never started because the active B build lock refused acquisition;
the general runner now holds an array until its locks and dependent aggregate
are established. Its 31 focused runner/oracle tests pass.

As of 18:41 UTC, corrected Tier C array \texttt{1723299} had 46/60 task
records (23 complete pairs): 30 SAT and 16 UNSAT, no recorded model/proof
failure or expected-verdict mismatch, and identical paired verdicts and proof
sizes. The remaining tasks and aggregate \texttt{1723300} are active; these
partial records are not an aggregate performance result. The checker bundle
is frozen by SHA-256 in the run manifest. No mutable checker clone is used.

The fixed ASan+UBSan public-action/online suite passed 10/10; a separately
rebuilt six-test API/online subset and the scheduling test passed. The latter
had a parent-identical stale witness: early Lucky refuted its old tiny binary
formula before the collection boundary under test. A deeper independent
UNSAT tail restores the test's intended path. \texttt{DefaultPolicy} and its
scheduler are now non-copyable/non-movable because their arms borrow
configuration or member addresses; this changes no solver schedule.

\section*{Learning design and open gates}
For inprocessing order/dose, the explicit online executor already has typed
steps, limits and replay, but the mature default scheduler separately owns
clocks, admission, cursors and follow-ups. A shared default-exact slot
controller is the first major refactor. For branching and clause deletion,
current sub-policy pointers are borrowed and cover only part of the decision;
owned typed arms, safe switch boundaries, state migration and complete
reduction-strategy control are the second. A third necessary seam is session
lifetime: \texttt{tactix\_solve} is synchronous and constructs a fresh
\texttt{DefaultPolicy} per call, so no live default-policy arm can yet be
reconfigured through the C API. Preserve default action/proof traces and
measure overhead before enabling a learner.

N2 is still open: the frozen canonical no-option DSR checker rejects the
separate H2 proof; the implied-unit deletion option is only a diagnostic
alternative with a different proof-state contract. Broad affected fuzz,
N1 completeness, N3/Tier-D competitiveness, N4 reference scheduling and the
retained Rank/B+C/WV gates remain open. The proof-rule audit found a qualified
winner gap in general live pseudo-Boolean arithmetic/redundancy/dominance;
ordinary CaDiCaL/Kissat clausal emitter events map to existing RUP/RAT/delete.
Finish Tier C and affected fuzz, settle and fully audit the H2 contract, then
run the cumulative gates. Only after those fixes should the 1,166-instance
Tier D pass use the frozen winner allocation: four CPUs, 32 GiB and 5,000
solver seconds per task on \texttt{cpu-dense-hwulab-q}.

\end{document}
